\chapter{Core grammar, records and execution details}
\label{app:contracts}

\section{Preserved Core 0 teaching grammar}

The grammar governing the preserved Core 0 teaching examples is reproduced
from \code{spec/core-0.lark}. It does not enlarge the first matched text
frontend beyond the portable profile described in the body. Names are
case-sensitive ASCII identifiers; strings are UTF-8 with JSON escaping. All
literal words in the grammar are reserved. The companion contract specifies
valid dates, exact rational literals, argument arity, type/context checks and
package resolution. A syntax parse is not semantic validation.

\lstinputlisting[basicstyle=\scriptsize\ttfamily]{spec/core-0.lark}

\section{Required core record definitions}

The historical closed machine-readable definitions remain in the original
design directory at \path{spec/core-ir.schema.json}. They define contexts, source spans,
observations, narratives, coverage units, mappings, partitions, entity/control
records, books, accounts, symbols, choices, expressions, constraints, assumptions,
alternatives, contracts, events, rules, packages, measurements, states,
journals, scenarios, information views, review decisions, evidence views,
model scopes, historical families, probability models, report mappings,
result cells, branches, runs and diagnostics. For those historical records the schema fixes field spelling and structural
shape, and its companion contract fixes reference resolution and operations.
New portable records instead obey the exact versioned contracts introduced in
Chapters~\ref{ch:language}, \ref{ch:execution} and \ref{ch:implementation}.
No historical field list grants new posting authority.

\begin{minipage}{\linewidth}
For example, these selected records reproduce the schema's field names.
Each field is required; a nullable link is represented by explicit null.

\begin{lstlisting}
EvidenceView {
  id, cutoff, system_as_of, mode,
  documents[], observations[], decisions[],
  conflict_sets[][], transforms[]
}
Context {
  id, entity, perimeter, concept, basis, measurement,
  time, unit, currency, taxonomy_version, dimensions[]
}
Symbol {
  id, type, context, domain, uncertainty_kind,
  scope, dependencies[], admission, provenance[]
}
Journal {
  id, event, book, date, accrual_period,
  postings[], rule, measurements[], assumptions[],
  precondition_hash, explanation
}
ResultCell {
  id, context, history_condition, scenario, path, kind,
  value, expression, lower, upper,
  lower_attained, upper_attained, sample_ref,
  certificate, reason, dependencies[]
}
\end{lstlisting}
\end{minipage}

A value is a reduced rational, or a typed symbolic/domain reference. A
posting is nonnegative with its side carrying the sign, in the book's
functional currency. Original units and the valuation reference survive
conversion. Monetary contexts require currency; a period flow requires a
calendar and exclusive end boundary. Historical zero postings may preserve traceability but establish no economic
event. New portable proposals omit zero lines and retain their explanation in
the event and findings instead. The amount may have
finer precision before the declared posting-rounding operation.

Schema validation rejects unknown fields but cannot prove that a treatment
is permitted, a reference points to the correct kind of record, or two periods
are comparable. Those semantic checks precede solver or runtime dispatch.
The structural fixture and complete demonstration source/package are supplied
separately so that a future implementation can check both layers explicitly.
\path{fixtures/core-ir-demo.json} supplies an authored expected lowering;
the reference checker evaluates its selected debt expressions. A general
source-to-IR compiler and complete external-registry validation remain to be
implemented. Domain records retain disclosed display text and rounding rules
alongside their exact bounds.

\section{Typing and partial-information rules}

For expressions $x$ and $y$ in the same currency $C$, the money addition rule is
\[
 \frac{x:\operatorname{Money}\langle C\rangle\quad
       y:\operatorname{Money}\langle C\rangle}
      {x+y:\operatorname{Money}\langle C\rangle}.
\]
Context compatibility is a separate premise for observations. There is no
analogous implicit rule for different currencies. Explicit conversion uses
a dated measurement and records its rate and purpose. A rate-times-principal
calculation is typed through a declared accrual convention; the example's
annual simple rate and year fraction produce money in the principal's currency.

Known expressions evaluate normally. Shared symbols evaluate to symbolic
expressions subject to a joint constraint store. Missing source values cannot
enter arithmetic until the model supplies a symbol or a conditional assumption.
Boolean guards over symbols produce formulas, not the value false. This avoids
the common error in which ``unknown whether impaired'' silently means
``not impaired.'' Division requires a proven nonzero denominator or a
partitioned domain that excludes zero; otherwise the expression remains
unsupported for concrete execution.

Logical queries first require a nonempty feasible family. An empty family
returns \code{Infeasible}; it must not produce vacuous answers to substantive
questions. For a nonempty family, constraint evaluation has three results.
\code{entailed} means every state in
the declared feasible family satisfies the proposition. \code{refuted} means
none does. \code{undetermined} means both remain feasible or the solver has
not decided; the response distinguishes those causes. Probabilistic truth is
a separate query requiring a measure. A language-level Boolean value must
never be inferred from an arbitrary probability threshold.

For a decided answer, let $F(h)$ be the exact constraint representation of the
whole family requested by the query, and let $P(h)$ be its proposition.
Establish feasibility first. Then \code{entailed} requires that
$F(h)\land\neg P(h)$ has no solution, while \code{refuted} requires that
$F(h)\land P(h)$ has no solution. Feasible witnesses on both sides establish
the mixed case. A timeout, unsupported operation, or unfinished search
establishes neither universal conclusion. A proof for a selected cell applies
to that cell; it applies to the parent family only with coverage of the
remaining cells or a proof over the full family. The result records the
scope, query kind and evidence for that status.
Section~\ref{sec:query-certainty} distinguishes this Boolean test from a
certain-tuples view and a family of answer schedules.


\section{Operational rule selection}

Let $V(e,h)$ be the applicable rules for event $e$ in concrete history $h$,
after package/date selection and guard evaluation. Let $r\prec r'$ mean that
$r'$ validly overrides $r$ through the declared acyclic graph. Define
\[
 \operatorname{Max}(V)=\{r\in V:\nexists r'\in V\text{ with }r\prec r'\}.
\]
For one required recognition operation, a singleton selects that rule; the
empty set yields \code{UnsupportedTreatment}; a larger set yields
\code{RuleConflict}. Multiple distinct operations, such as measurement and
recognition, are evaluated through separate declared interfaces. A rule that
intentionally produces no postings must belong to a supported non-recognition
event operation, whose economic state effect is still recorded.

Given a selected rule, the evaluator computes all local bindings and
preconditions, emits proposed postings, verifies
Equation~\eqref{eq:posting}, validates account/context restrictions and commits
the event transaction. Failure returns the original state plus a diagnostic.
This defines the atomicity requirement independently of the storage technology.

\begin{minipage}{\linewidth}
\section{Example run request and response}

The following request executes two named conditions under one scenario. It
does not request an aggregate probability distribution.

\begin{lstlisting}
{
  "model": "BorrowingExample@1",
  "evidence_view": "synthetic-2025-exact",
  "starting_state": "after-2025-close",
  "horizon": {"start": "2026-01-01", "end_exclusive": "2027-01-01"},
  "conditions": ["allFixed", "allFloating"],
  "scenarios": ["rateShock"],
  "mode": "conditional",
  "weight_model": null,
  "reports": ["balance_sheet", "income", "cash_flow", "equity", "journals"],
  "package_lock": "demonstration-packages-v1"
}
\end{lstlisting}
\end{minipage}

The response contains two separately labelled results with interest expense
5 and 7, cash 90 and 88, debt 100 in both, and equity $-10$ and $-12$.
A request for the parent symbolic family returns the formula
$7-2\alpha$ and, under exact linear bounds, $[5,7]$ before monetary
quantisation. The API rejects an unconditional mean request without a weight
model. A user can instead request an arbitrary clearly labelled representative
case, which does not stand in for the feasible family.

\section{Requirement traceability}

\begin{center}
\begin{tabular}{lll}
\toprule Requirement & Primary implementation area & Conformance cases\\\midrule
R1 & Evidence registry and extraction adapters & T03, T09\\
R2 & Context/version model and cutoff queries & T03, T04\\
R3 & Types, mappings and package selection & T02, T12, T13\\
R4 & Symbols, choices and constraint graph & T05--T10, T14\\
R5 & Posting, closing and consolidation & T01, T11--T13\\
R6 & Contract and behavioural interfaces & T05, T11, T15, T16\\
R7 & Conditional output and probability types & T06, T07, T18\\
R8 & Explanation graph and report dependencies & T09, T11--T13\\
R9 & Capabilities and versioned modules & T02, T14--T16\\
R10 & Run records, scheduler and replay & T03, T17, T18\\\bottomrule
\end{tabular}
\end{center}
